\subsection{Materials for tagless-protein purification}

Protein Select resin (100~mL; product 17542102) for purification of tagless ZAPP-1 was purchased from Cytiva.

Amicon Ultra-15 centrifugal filters (15~mL capacity; MilliporeSigma) with 3~kDa (UFC900396) and 10~kDa (UFC901096) MWCO membranes were obtained from Fisher Scientific. Size-exclusion chromatography was performed on an \"{A}KTA pure chromatography system using Superdex 200 Increase 10/300 GL and Superdex 75 Increase 10/300 GL columns, all purchased from Cytiva.

\subsection{Tagless-protein production and purification for characterization}

To assess whether the C-terminal Strep-tag II affected ZAPP-1 activity, we produced tagless ZAPP-1 using the split-intein-based Cytiva Protein Select purification system \citep{clifford2024production}. This approach allows flexibility in buffer composition, with tag cleavage and elution performed without added protease, imidazole, or biotin. ZAPP-1 was cloned by BsaI-mediated Golden Gate assembly into a custom T7 expression vector encoding an N-terminal His$_6$--Protein Select fusion in the format MSHHHHHHSG-MVKIVSRKSLGVQNVYDIGVEKDHNFLLANGLIASN-SVG-design-GS. On-column intein cleavage removed the N-terminal fusion and yielded SVG-design-GS. Here, design denotes the 198-residue ZAPP-1 design sequence. Expression followed the large-scale protocol described above at a 600~mL culture scale.

Frozen cell pellets were thawed and processed using a similar lysis and clarification protocol to that described in the preceding section, with 20~mL of lysis buffer per 50~mL of culture (240~mL total). The base purification buffer consisted of 40~mM HEPES, 50~mM NaCl, and 200~$\mu$M ZnSO$_4$, pH~8.0. The corresponding lysis buffer contained the same supplements as described above, and washing, cleavage, and elution were all performed in the base purification buffer. Buffers were prepared and filtered as described above.

Clarified lysate was divided equally between two gravity-flow columns containing 12~mL beds of Cytiva Protein Select resin (120~mL lysate per column). The combined flow-through was applied to an additional $\sim$20~mL resin bed to recover unbound protein. Each column was washed with $3 \times 3.5$ column volumes (CV) of base purification buffer. The column outlets were then closed, and $\sim$3~mL of base purification buffer was added above each resin bed to maintain a layer of buffer over the resin during incubation. Columns were incubated at room temperature for 20--24~hrs to allow intein-mediated tag cleavage. The outlets were then opened, and tagless protein was eluted with $3 \times 1$~CV of the same buffer.

Collected eluates were concentrated to $\sim$6.6~mL. Insoluble material that formed during concentration was removed by centrifugation, and the soluble protein was purified by size-exclusion chromatography on a Superdex 75 Increase 10/300 GL column using the base purification buffer. Monomeric fractions were supplemented with NaHCO$_3$ to $\sim$45~mM using a 500~mM stock prepared in Milli-Q water and incubated at room temperature with shaking for $\sim$2~hrs. Protein was then concentrated using an Amicon Ultra-15 centrifugal concentrator (10~kDa MWCO) at 2000$g$ to $\sim$1.4~mL ($\sim$60~mg/mL). Protein molecular weight and tag removal were confirmed by intact-protein ESI-MS (\cref{fig:PTE_v2SI_ZAPP1_tagless}).

Used Protein Select resin was regenerated with $3 \times 5$~CV of 30\% (v/v) isopropanol containing 100~mM NaOH, with resin resuspension to ensure thorough mixing and a contact time of $>15$~min. The resin was then washed with $3 \times 5$~CV of Milli-Q water and stored in 40\% (v/v) ethanol in water. Elution directly into a protein-compatible buffer, without added protease, imidazole, or biotin, may simplify preparative-scale production of tagless protein for biochemical characterization and crystallization screening.
\subsection{Tagless-protein kinetic analysis}

Tagless ZAPP-1 was assayed under the same bicarbonate-supplemented conditions at \([E]_0 = 23.4\)~$\mu$M, with a 30~min acquisition and initial velocities fit from 300 to 1800~s. Measurements were performed in technical triplicate at each substrate concentration.

For the tagged-versus-tagless comparison, approximate pointwise 95\% confidence intervals for the fitted Michaelis-Menten curves were calculated from the full covariance matrix of the unweighted nonlinear regression using the delta method and a Student's $t$ multiplier with four degrees of freedom (six substrate-concentration means minus two fitted parameters). These intervals describe uncertainty in the fitted mean response and include regression uncertainty only, excluding enzyme-concentration and calibration uncertainties.
